% Shared preamble for the three Weilight Harbor documents.
% Compile each main file with pdfLaTeX (Overleaf default).

\documentclass[11pt,a4paper]{article}

% --- Page layout
\usepackage[a4paper,margin=2.4cm]{geometry}
\usepackage{setspace}
\setstretch{1.10}
\setlength{\parskip}{4pt}
\setlength{\parindent}{0pt}

% --- Fonts and basics
\usepackage[utf8]{inputenc}
\usepackage[T1]{fontenc}
\usepackage{microtype}
\usepackage{lmodern}
\usepackage{textcomp}
% Map Unicode em/en-dashes and the simple bullet to LaTeX equivalents so they
% render in the default Computer Modern font without falling back to question
% marks.
\DeclareUnicodeCharacter{2014}{---}
\DeclareUnicodeCharacter{2013}{--}
\DeclareUnicodeCharacter{2026}{\ldots}
\DeclareUnicodeCharacter{2018}{`}
\DeclareUnicodeCharacter{2019}{'}
\DeclareUnicodeCharacter{201C}{``}
\DeclareUnicodeCharacter{201D}{''}
\DeclareUnicodeCharacter{2022}{\textbullet}

% --- Common packages
\usepackage{graphicx}
\usepackage{tabularx}
\usepackage{booktabs}
\usepackage{array}
\usepackage{longtable}
\usepackage{enumitem}
\setlist{nosep,leftmargin=*}
\usepackage{xcolor}
\usepackage{hyperref}
\hypersetup{
  colorlinks=true,
  linkcolor=black,
  citecolor=black,
  urlcolor=blue!55!black,
  pdfborder={0 0 0}
}
\usepackage{titlesec}
\titleformat{\section}{\large\bfseries}{\thesection}{0.6em}{}
\titleformat{\subsection}{\normalsize\bfseries}{\thesubsection}{0.5em}{}
\titleformat{\subsubsection}{\normalsize\itshape}{\thesubsubsection}{0.5em}{}
\titlespacing*{\section}{0pt}{1.6ex}{0.8ex}
\titlespacing*{\subsection}{0pt}{1.2ex}{0.5ex}
\titlespacing*{\subsubsection}{0pt}{1.0ex}{0.3ex}

% Inline code
\usepackage{listings}
\lstset{
  basicstyle=\ttfamily\footnotesize,
  breaklines=true,
  columns=fullflexible,
  showstringspaces=false,
  frame=single,
  framerule=0.3pt,
  rulecolor=\color{black!30},
  xleftmargin=0pt,
  xrightmargin=0pt,
  aboveskip=4pt,
  belowskip=4pt
}

% Tighter table column types
\newcolumntype{L}[1]{>{\raggedright\arraybackslash}p{#1}}
\renewcommand{\arraystretch}{1.15}

% Header / footer
\usepackage{fancyhdr}
\pagestyle{fancy}
\fancyhf{}
\fancyfoot[C]{\thepage}
\fancyhead[L]{\small Weilight Harbor — Group 14}
\fancyhead[R]{\small COMP3030J 2025/26}
\renewcommand{\headrulewidth}{0.3pt}


\title{\vspace{-1.5cm}\textbf{Weilight Harbor — System Document}\\
       \large COMP3030J Web Development, UCD, 2025/26 — Group 14 (Draft, Week 11)}
\author{}
\date{}

\begin{document}
\maketitle
\vspace{-2.0cm}
\begin{center}
  \small
  Live deployment: \url{http://csi6220-2-vm3.ucd.ie} \quad
  Repository: \url{https://github.com/licctvcctv/weilight-harbor}
\end{center}
\vspace{0.4cm}

% =====================================================================
\section*{Abstract}
\addcontentsline{toc}{section}{Abstract}

Weilight Harbor (\textit{Weiguang Gangwan}, ``harbor of faint lights'') is a Flask-based
web platform built for the families and caregivers of seriously ill patients.
The system unifies five modules that mainstream platforms typically keep
separate: a moderated peer community, a time-gated anonymous \emph{Night
Confessions} space, a verified crowdfunding flow with simulated donations,
a private journal with an \emph{Annual Film} reflection, and a Leaflet-based
\emph{Respite Station} map for arranging temporary care help. Three roles
are supported via RBAC inherited from the Pear-Admin-Flask base: regular
visitor, certified family/volunteer, and administrator. The application is
served by Gunicorn behind Nginx on an Ubuntu VM provided by UCD, persists
to a single SQLite database with twelve tables, and supports full
Mandarin/English internationalisation through Flask-Babel
(\textbf{511 translated strings}). Continuous deployment is automated by a
GitHub Actions workflow that traverses the UCD jump host with
\texttt{expect} and a TOTP secret, force-resets the working tree to
\texttt{origin/main}, and reloads Gunicorn. This document records the
design rationale, the module-by-module implementation, the use of
generative AI in the build, and the per-member contribution split.

% =====================================================================
\section{Introduction}

\subsection{Background and motivation}
Caregivers of seriously ill patients face emotional exhaustion, social
isolation, and financial strain that mainstream platforms address only in
fragments. Generic forums lack verification; fundraising sites are
purely transactional; chat groups are ephemeral and private. Weilight
Harbor was designed as a \emph{quiet, warm} space where these users can
both be \emph{heard} and \emph{helped} — emotional companionship before
financial transaction, and verified identity before public appeal.

\subsection{Goals}
\begin{itemize}
  \item Provide a verified community where certified caregivers and the
        public can interact with low friction and clear safety signals.
  \item Enable trustworthy, simulated crowdfunding to demonstrate the
        end-to-end donation flow, without handling real money.
  \item Offer private journaling with mood tracking and a year-end
        \emph{Annual Film} reflection.
  \item Surface practical mutual aid (respite and equipment lending) on
        an interactive map.
  \item Support both Mandarin and English audiences with one codebase.
\end{itemize}

\subsection{Stakeholders and roles}
\begin{itemize}
  \item \textbf{Regular user} — read, donate, post in the community,
        comment, react.
  \item \textbf{Certified family / volunteer} — additionally create
        crowdfunding campaigns, accept respite requests, and access the
        private journal.
  \item \textbf{Administrator} — additionally review certification and
        campaign applications, and moderate community content.
\end{itemize}

\subsection{Document structure}
Section~\ref{sec:tech} describes the technical implementation, end-to-end.
Section~\ref{sec:ai} evaluates how generative AI tooling was used during
development. Section~\ref{sec:team} records team contribution.
Section~\ref{sec:concl} concludes. Appendices list the routing surface,
seeded test accounts, and a sample LCP/Lighthouse run.

% =====================================================================
\section{Technical Implementation}\label{sec:tech}

\subsection{Technology stack}
\begin{tabularx}{\linewidth}{@{} l X @{}}
  \toprule
  \textbf{Layer} & \textbf{Choice and rationale} \\
  \midrule
  Backend         & Python 3.10 with Flask 3.x, application-factory pattern. \\
  ORM             & SQLAlchemy + Flask-Migrate; Alembic-style migrations. \\
  Auth            & Flask-Login session cookies; Werkzeug PBKDF2 hashing. \\
  i18n            & Flask-Babel with a single \texttt{messages.po}
                    catalogue (511 strings translated). \\
  Frontend        & Server-rendered Jinja2 templates with Bootstrap 5.3,
                    Lucide icons, and Leaflet.js for the map. Minimal
                    custom JS. \\
  Database        & SQLite (single-file), tuned for the assignment scale. \\
  Process         & Gunicorn 21 sync workers (default 2). \\
  Reverse proxy   & Nginx 1.18 — TLS termination point and static asset
                    server. \\
  Service         & systemd, with Gunicorn behind a managed socket. \\
  CI/CD           & GitHub Actions + \texttt{expect} script through the
                    UCD jump host. \\
  Admin base      & Pear-Admin-Flask (Layui) — reused for RBAC, dict,
                    and audit log scaffolding. \\
  \bottomrule
\end{tabularx}

\subsection{High-level architecture}
\begin{verbatim}
            Browser
              | HTTP
              v
       Nginx :80  -- /static/ served directly
              |
              v (Unix socket)
       Gunicorn (WSGI, sync workers, systemd-managed)
              |
              v
       Flask app (create_app factory)
        - blueprints: auth, community, crowdfunding,
          journal, respite, user_center, index, admin
        - Flask-Login, Flask-Babel, Flask-WTF, Flask-Reuploaded
              |
              v
       SQLAlchemy
              |
              v
       SQLite (12 tables, ~100 KB seeded)
\end{verbatim}

\subsection{Source layout}
\begin{lstlisting}
weilight-harbor/
  app.py                        # WSGI entry point
  applications/
    __init__.py                 # create_app factory
    configs/                    # SQLite, Babel, upload paths
    extensions/                 # init_sqlalchemy/init_login/init_babel...
    models/                     # SQLAlchemy models (12 tables)
    view/                       # Route blueprints
      auth/  community/  crowdfunding/
      journal/  respite/  user_center/
      index/  admin/
    common/utils/               # files, http, sensitive, rights, ...
  templates/
    public/                     # 31 Bootstrap 5 templates
    admin/                      # Layui admin panel
  static/
    public/  admin/
  translations/zh/LC_MESSAGES/  # 511 strings
\end{lstlisting}

\subsection{Data model}
The schema is built from twelve SQLAlchemy models. The user/role/power
triplet is inherited from Pear-Admin-Flask and reused as-is. The
\texttt{wl\_*} tables are domain-specific.

\begin{tabularx}{\linewidth}{@{} L{3.4cm} X @{}}
  \toprule
  \textbf{Table} & \textbf{Notable columns} \\
  \midrule
  \texttt{admin\_user}
    & \texttt{username, email, phone, password\_hash, user\_type,
       is\_certified, locale, avatar, bio} \\
  \texttt{admin\_role}, \texttt{admin\_user\_role},
  \texttt{admin\_role\_power}, \texttt{admin\_power}
    & RBAC mapping inherited from the base. \\
  \texttt{wl\_post}
    & \texttt{user\_id, title, content, category, images} (JSON),
      \texttt{like\_count, comment\_count, view\_count, is\_anonymous,
      delete\_at}. \\
  \texttt{wl\_comment}
    & \texttt{post\_id, user\_id, content, parent\_id} (one-level
      threading). \\
  \texttt{wl\_post\_like}
    & \texttt{(post\_id, user\_id) UNIQUE} for idempotent likes. \\
  \texttt{wl\_confession}
    & \texttt{content, nickname, session\_id, hug\_count} (Night
      Confessions). \\
  \texttt{wl\_campaign}
    & \texttt{user\_id, title, description, funding\_goal,
      current\_amount, status} ($\in \{$pending, active, closed,
      rejected$\}$), images. \\
  \texttt{wl\_donation}
    & \texttt{campaign\_id, user\_id, amount, message, is\_anonymous}. \\
  \texttt{wl\_journal\_entry}
    & \texttt{user\_id, title, content, images} (JSON),
      \texttt{mood, entry\_date}. \\
  \texttt{wl\_respite\_request}
    & \texttt{user\_id, request\_type} (service/equipment),
      \texttt{lat, lng, status, acceptor\_id}. \\
  \texttt{wl\_certification}
    & \texttt{user\_id, cert\_type, real\_name, id\_card,
      document\_url, additional\_docs} (JSON), \texttt{status}. \\
  \texttt{wl\_sensitive\_word}
    & \texttt{word, severity} ($1$=warn, $2$=crisis modal). \\
  \bottomrule
\end{tabularx}

\smallskip
All domain tables inherit a common timestamp triple
(\texttt{create\_at}, \texttt{update\_at}, \texttt{delete\_at}) for
soft-delete and audit; \texttt{delete\_at IS NULL} is the universal
"alive" predicate.

\subsection{Authentication and authorisation}
\textbf{Authentication.} Flask-Login session cookies; passwords stored as
PBKDF2-SHA256 hashes via \texttt{werkzeug.security}. Login accepts
\emph{username, email, or phone} in the same field; a uniqueness
constraint at the model level prevents collisions during registration.
Logout is POST-only to defeat image-tag CSRF.

\textbf{Authorisation.} Public-side blueprints use the standard
\texttt{@login\_required} decorator (\textgreater 60 protected endpoints
in the codebase). Admin-side endpoints add a custom
\texttt{@authorize("admin:\textit{module}:\textit{action}", log=True)}
decorator that resolves the current user's effective power codes through
\texttt{admin\_user $\to$ admin\_user\_role $\to$ admin\_role $\to$
admin\_role\_power $\to$ admin\_power} and writes an audit row for every
permitted action. Negative cases short-circuit with HTTP 403 (admin) or
302 to \texttt{/auth/login?next=…} (public).

\subsection{Module-by-module walkthrough}

\subsubsection{Auth (\texttt{/auth/*})}
\begin{itemize}
  \item \texttt{login}: identifier-agnostic — same field accepts username,
        email or phone; \texttt{User.validate\_password} drives the
        Werkzeug check.
  \item \texttt{register}: minimum length (6), uniqueness on email/phone;
        new accounts default to the \emph{regular} role.
  \item \texttt{forgot-password}: UI is wired; SMTP is stubbed because
        the assigned VM has no outbound mail credentials.
  \item \texttt{logout}: POST-only.
\end{itemize}

\subsubsection{Community (\texttt{/community/*})}
The community page uses a 70/30 desktop layout: feed on the left, a
sticky \emph{Night Confessions} panel on the right. Posts support
anonymous mode (renderer hides username and avatar). Liking is AJAX, with
the underlying \texttt{(post\_id, user\_id)} unique index providing
idempotent toggles. Comments support \texttt{parent\_id} threading.
Sensitive-word checking is two-tier: severity~1 raises a soft warning;
severity~2 surfaces a crisis hotline modal client-side and short-circuits
submission server-side. \emph{Night Confessions} is gated by server-local
clock to the 21:00–06:00 window; outside the window the panel renders a
countdown.

\subsubsection{Crowdfunding (\texttt{/crowdfunding/*})}
Campaigns can be created only by certified users; an admin must approve
the campaign before it appears in the public listing. A donation is a
single \texttt{POST} that increments \texttt{current\_amount} inside an
SQLAlchemy transaction; the response payload includes the new percentage
so the front-end progress bar can animate. A campaign auto-flips to
\emph{Fully Funded} when \texttt{current\_amount $\geq$ funding\_goal};
the donate button is greyed out and the badge appears on every card.

\subsubsection{Life Journal (\texttt{/journal/*})}
Every query is filtered by \texttt{user\_id == current\_user.id}; no
operator role can read a user's journal through the public site.
Annual Film selects up to twelve entries from a chosen year, ordered by
month, and renders a full-screen Ken-Burns slideshow on the client.

\subsubsection{Respite Station (\texttt{/respite/*})}
Backed by Leaflet on top of OpenStreetMap tiles. Requests cluster via
\texttt{Leaflet.markercluster}; orange pins are service requests, green
pins are equipment lending. The state machine is
\texttt{pending $\to$ in\_progress $\to$ completed}. A request cannot be
accepted by its creator. Mobile layout swaps the side panel for a bottom
sheet.

\subsubsection{User Centre (\texttt{/user/*})}
Profile, edit, settings, certification application, certification
status, and a paginated history page with four tabs (Posts, Campaigns,
Donations, Respite Requests). Avatar uploads use Flask-Reuploaded with an
extension allow-list. Pagination links are constructed via
\texttt{url\_for(\textit{endpoint}, **\{page\_param: n\})} so query
strings never contain the brittle ad-hoc \texttt{\&} concatenation that
the early prototype used.

\subsubsection{Admin (\texttt{/admin/*})}
Three custom blueprints sit on top of the Pear-Admin-Flask base:
\texttt{adminCert} (certification queue), \texttt{adminCampaign}
(campaign queue), and \texttt{adminCommunity} (post moderation). Every
state-changing action is wrapped in \texttt{@authorize(..., log=True)}
which writes a row to the audit log. Reviewers can approve, reject (with
required reason), or hide; deletes are soft.

\subsection{Internationalisation}
All user-visible strings are wrapped with \texttt{\{\{ \_('…') \}\}}
(Jinja) or \texttt{gettext(…)} (Python). The catalogue lives at
\texttt{translations/zh/LC\_MESSAGES/messages.po} and contains 511
translated strings. Locale resolution order is: query string
(\texttt{?locale=}) $\to$ session $\to$ \texttt{Accept-Language}. The
\textbf{ZH/EN} toggle in the navbar persists the choice via both session
and \texttt{localStorage}, so a returning visitor on the same browser
sees their last choice.

\subsection{Front-end design system}
The platform uses CSS variables for a warm "harbor" palette
(\texttt{--color-primary}, \texttt{--color-primary-bg},
\texttt{--color-text-primary}, \texttt{--color-bg-card},
\texttt{--shadow-sm}, \texttt{--radius-md},
\texttt{--transition-fast}, etc.). Bootstrap 5 utility classes drive
layout; component skins (\texttt{wl-tag}, \texttt{wl-progress},
\texttt{btn-wl-primary}) wrap Bootstrap defaults. Lucide icons are
rendered through \texttt{data-lucide} and initialised once per page on
\texttt{DOMContentLoaded}.

\subsection{Validation, sanitisation and uploads}
All form fields are length-clamped (\texttt{[:N]}) and re-validated on
the server side regardless of client behaviour. Email and phone are
uniqueness-checked. File uploads pass through a single helper
\texttt{save\_upload(file, folder, prefix, allowed\_exts=…)} that:
\begin{enumerate}
  \item verifies the extension against an explicit allow-list;
  \item generates a UUID-prefixed filename so collisions and path
        traversal are impossible;
  \item writes to \texttt{static/uploads/<folder>/}.
\end{enumerate}
Sensitive-word checking is debounced 300\,ms client-side; the
authoritative decision still runs on the server.

\subsection{Deployment}
Two deploy paths are supported, sharing the same \texttt{vm\_update.sh}
on the VM.

\textbf{Local one-click (\texttt{deploy.sh}).}
The script generates a TOTP for the UCD jump host, drives \texttt{expect}
through \texttt{jumper $\to$ VM}, runs
\texttt{git fetch origin main \&\& git checkout -fB main FETCH\_HEAD
\&\& git reset --hard FETCH\_HEAD \&\& git clean -fd \&\& bash
vm\_update.sh}. After a successful run, it captures the new
\texttt{HEAD} on the VM and asserts it matches the local branch SHA.
A diagnostics-only mode (\texttt{./deploy.sh --diag}) and a vendor sync
(\texttt{./deploy.sh --sync-vendor}) are exposed for the gitignored
heavy assets (Layui, TinyMCE, ECharts).

\textbf{GitHub Actions (\texttt{.github/workflows/deploy.yml}).}
Triggered on \texttt{push} to \texttt{main}; uses the same
\texttt{expect} pattern with secrets held in \emph{GitHub Secrets}.
Concurrency is locked to a single deployment at a time.

\subsection{Security considerations}
\begin{itemize}
  \item Sessions are HTTPOnly cookies; passwords PBKDF2-SHA256.
  \item Write actions are POST and CSRF-protected via Flask-WTF where
        the form is rendered through the helper.
  \item File uploads validated by extension and isolated under
        \texttt{static/uploads/}.
  \item Admin endpoints behind RBAC + audit log.
  \item \texttt{.gitignore} excludes \texttt{*.db}, \texttt{*.env},
        large vendor JS files, and \texttt{.DS\_Store}.
  \item \textbf{Limitations} (acknowledged): no rate limiting; no
        CAPTCHA; the sensitive-word list is finite and bypassable; the
        platform is intentionally \emph{not} PCI-compliant — donations
        are simulated.
\end{itemize}

\subsection{Testing strategy}
We use a manual smoke-test matrix covering the primary user journeys for
each role, driven by the eight seeded demo accounts (Section~A.2). The
\texttt{flask seed} and \texttt{flask seed-demo} CLI commands produce a
reproducible state (8 users, 12 community posts, 7 campaigns, 82
donations, 50 journal entries, 8 respite requests, 15 confessions). A
small per-module Python smoke probe (used during pull-request review)
hits each blueprint's index route and asserts that the response is
either 200 or a 302 redirect to login.

\subsection{Performance and accessibility}
A baseline Lighthouse run on the live VM produces, on the landing page,
\textbf{Performance $\geq$ 90} and \textbf{Best Practices $\geq$ 95} for
desktop. Server-rendered pages are below 25\,KB gzipped; static assets
are served directly by Nginx with a 30-day cache. WCAG AA has \emph{not}
been formally validated; partial measures include semantic HTML,
keyboard-reachable interactive controls, and visible focus rings on
custom components.

\subsection{Limitations and known issues}
\begin{itemize}
  \item SQLite write contention under heavy concurrent load (acceptable
        for the assignment scale).
  \item Forgot-password flow is UI-only — no real e-mail is sent.
  \item Annual Film unlocks only if a user has $\geq$10 entries in a year.
  \item Public OpenStreetMap tile rate limits apply on Respite Station.
\end{itemize}

% =====================================================================
\section{Generative AI Use}\label{sec:ai}

\subsection{Tools used}
\begin{itemize}
  \item \textbf{ChatGPT (GPT-4 / GPT-4o)} — architectural reasoning,
        microcopy drafting, debugging.
  \item \textbf{Claude (Sonnet / Opus)} — long-context refactors across
        multiple files, Jinja template clean-up, deployment-script
        reasoning.
  \item \textbf{GitHub Copilot} — line-level autocompletion in VS Code.
\end{itemize}

\subsection{Where AI helped}
\begin{tabularx}{\linewidth}{@{} L{4cm} X @{}}
  \toprule
  \textbf{Area} & \textbf{Use} \\
  \midrule
  Initial scaffolding
    & Drafting blueprint folder structure, route stubs, model fields. \\
  Bilingual microcopy
    & First-pass Mandarin/English for empty states, error messages,
      modals; team revised tone to fit the warm/quiet palette. \\
  Refactors
    & Extracting shared decorators, normalising form-validation,
      cleaning Jinja macros (e.g.\ the \texttt{render\_history\_pagination}
      macro on the user history page). \\
  Deployment debugging
    & Diagnosing why the GitHub Actions deploy reported
      \texttt{\_\_RC\_\_:0} while the live site never changed (root
      cause: missing \texttt{cd \~/weilight-harbor} in the remote
      command); diagnosing why the VM repo was stuck on an empty
      \texttt{master} branch and never advanced past the initial
      \texttt{deploy.sh --tar} import. \\
  Documentation
    & First drafts of \texttt{README.md}, this system document, and
      the user document. \\
  \bottomrule
\end{tabularx}

\subsection{Where AI was avoided}
\begin{itemize}
  \item Decisions on data model, role boundaries, and module split were
        taken by the team based on the product brief.
  \item Sensitive-word list, illness categories, and crisis-hotline copy
        were authored manually to avoid hallucinated or insensitive
        content.
  \item All AI-suggested code was opened, read top-to-bottom, and
        manually tested before commit.
\end{itemize}

\subsection{Verification practices}
\begin{enumerate}
  \item AI-suggested files were reviewed line by line and run on a fresh
        SQLite database before being merged.
  \item AI-drafted SQL or migration logic was re-checked against
        \texttt{db.create\_all()} on a clean database file.
  \item Deployment scripts were dry-run with the diagnostics-only mode
        (\texttt{./deploy.sh --diag}) before any change touched the
        production VM.
\end{enumerate}

\subsection{Reflection}
Generative AI most clearly accelerated \emph{plumbing}: form rendering,
i18n wrapping, repeated validation, and shell scripts. It was less
useful for \emph{taste-sensitive} work — palette, microcopy tone,
illness terminology — where unchecked output would have damaged the
empathic posture the platform is built around. Net effect:
approximately $30$–$40\%$ effort saved on plumbing, with no measurable
saving on design judgement.

% =====================================================================
\section{Team Contribution}\label{sec:team}

\textit{The contribution table below is a placeholder for the Week~11
draft. Final names and percentages will be confirmed before the Week~13
final submission. Each row is intended to be replaced with: full name,
student number, modules owned, key commits, and \% contribution.}

\subsection{Contribution matrix (placeholder)}
\begin{tabularx}{\linewidth}{@{} L{3.0cm} L{4.5cm} X r @{}}
  \toprule
  \textbf{Member} & \textbf{Modules / areas} & \textbf{Notable deliverables} & \textbf{\%} \\
  \midrule
  A (TBD) & Auth, User Centre, Certification & Login/Register, Profile,
            Certification flow, History page & TBD \\
  B (TBD) & Community, Sensitive Words, Night Confessions
          & Feed UI, like/comment, anonymous post, time-gated panel &
            TBD \\
  C (TBD) & Crowdfunding, Donation Flow
          & Campaign listing/detail, simulated donate, progress
            animation & TBD \\
  D (TBD) & Journal, Annual Film & Timeline UI, mood tracking,
            Ken-Burns slideshow & TBD \\
  E (TBD) & Respite Station, Map & Leaflet integration, marker
            clustering, accept/complete & TBD \\
  F (TBD) & Admin Panel, RBAC, Deployment, i18n
          & Admin review pages, Babel catalogue, Nginx/Gunicorn,
            GitHub Actions & TBD \\
  \bottomrule
\end{tabularx}

\subsection{Working practices}
\begin{itemize}
  \item Source control on GitHub with feature branches; \texttt{main} is
        deployed automatically on push.
  \item Issues tracked on the GitHub Issues board, grouped by module.
  \item Weekly synchronous meeting; asynchronous coordination on
        WeChat.
  \item Code review: at least one peer review on non-trivial PRs;
        AI-generated code requires the author to walk a peer through the
        diff before merge.
\end{itemize}

\subsection{Individual reflections (placeholder)}
\textit{Each member contributes a $\approx$150-word reflection here
before final submission.}

% =====================================================================
\section{Conclusion}\label{sec:concl}

Weilight Harbor delivers, end-to-end, a multi-module web platform that
places caregiver wellbeing at its centre. The five user-facing modules,
a back-office admin moderation layer, RBAC, full Mandarin/English
internationalisation, and an automated deployment pipeline are all live
on the assigned UCD VM. Generative AI tooling shortened plumbing work
but was deliberately kept out of taste-sensitive content. The known
limitations (SQLite scale, simulated donations, stubbed e-mail) are
acknowledged and documented; none are blockers for the academic scope
of the assignment. The remaining iteration before submission will focus
on user-acceptance feedback from Week~12 peer testing, accessibility
refinement, and final per-member reflection write-ups.

% =====================================================================
\appendix
\section*{Appendix A — Selected routes}
\addcontentsline{toc}{section}{Appendix A — Selected routes}
{\small
\begin{tabularx}{\linewidth}{@{} L{1.6cm} L{5.5cm} X @{}}
  \toprule
  \textbf{Method} & \textbf{Path} & \textbf{Purpose} \\
  \midrule
  GET/POST & \texttt{/auth/login} & Login \\
  GET/POST & \texttt{/auth/register} & Register \\
  GET      & \texttt{/community/} & Post feed \\
  GET      & \texttt{/community/post/<id>} & Post detail \\
  POST     & \texttt{/community/post/<id>/like} & Toggle like \\
  POST     & \texttt{/community/post/<id>/comment} & Add comment \\
  GET/POST & \texttt{/community/confessions} & Night Confessions
            (read/submit) \\
  GET      & \texttt{/crowdfunding/} & Campaign listing \\
  GET      & \texttt{/crowdfunding/campaign/<id>} & Campaign detail \\
  POST     & \texttt{/crowdfunding/campaign/<id>/donate} & Simulated
            donation \\
  GET      & \texttt{/journal/} & Private timeline \\
  GET/POST & \texttt{/journal/create} & New entry \\
  GET      & \texttt{/journal/annual-film} & Annual Film slideshow \\
  GET      & \texttt{/respite/} & Map page \\
  POST     & \texttt{/respite/create} & Create request \\
  POST     & \texttt{/respite/request/<id>/accept} & Accept request \\
  GET      & \texttt{/user/profile} & Personal centre \\
  GET      & \texttt{/user/history} & Activity history (paginated) \\
  GET/POST & \texttt{/user/certification} & Certification application \\
  GET      & \texttt{/admin/certification/} & Certification queue \\
  PUT      & \texttt{/admin/certification/approve/<id>} & Approve cert \\
  GET      & \texttt{/admin/campaign-review/} & Campaign queue \\
  GET      & \texttt{/admin/community/} & Community moderation \\
  POST     & \texttt{/set-locale} & Switch language \\
  \bottomrule
\end{tabularx}}

\section*{Appendix B — Seeded test accounts}
\addcontentsline{toc}{section}{Appendix B — Seeded test accounts}
{\small
\begin{tabularx}{\linewidth}{@{} l l L{3.4cm} X @{}}
  \toprule
  \textbf{Username} & \textbf{Password} & \textbf{Role} & \textbf{Notes} \\
  \midrule
  \texttt{admin}     & \texttt{admin123} & Administrator     & Full admin access \\
  \texttt{limei}     & \texttt{test123}  & Certified Family   & ALS caregiver, has campaigns \\
  \texttt{zhangwei}  & \texttt{test123}  & Regular User       & Donor profile \\
  \texttt{wangfang}  & \texttt{test123}  & Certified Volunteer & Retired nurse \\
  \texttt{chenyun}   & \texttt{test123}  & Certified Family   & Alzheimer's caregiver \\
  \texttt{liuyang}   & \texttt{test123}  & Regular User       & Pending certification \\
  \texttt{sunli}     & \texttt{test123}  & Certified Family   & Pediatric cancer caregiver \\
  \texttt{zhaojing}  & \texttt{test123}  & Certified Volunteer & Social worker \\
  \bottomrule
\end{tabularx}}

\section*{Appendix C — Codebase metrics}
\addcontentsline{toc}{section}{Appendix C — Codebase metrics}
\begin{tabularx}{\linewidth}{@{} l X @{}}
  \toprule
  Public route handlers (\texttt{@*\_bp.route|.get|.post}) & 63 \\
  Public-side HTML templates  & 31 \\
  Public-side JS modules      & 6 \\
  i18n catalogue (\texttt{messages.po}) lines & 2{,}582 \\
  Translated strings           & 511 \\
  Database tables              & 12 \\
  Seeded users / posts / campaigns / donations / journal / respite
                               & 8 / 12 / 7 / 82 / 50 / 8 \\
  \bottomrule
\end{tabularx}

\end{document}
